\chapter{Psicología del fundador: rasgos, hábitos y cognición experta}\label{ch:founder-psychology}

% Alcance: la psicología BASADA EN EVIDENCIA del emprendedor --- autoeficacia (Bandura +
%   metaanálisis de Stajkovic y Luthans), establecimiento de metas (Locke y Latham), rasgos
%   de personalidad vinculados a la tarea emprendedora (Rauch y Frese), resultados vs.\ inversiones
%   de capital humano (Unger et al.), práctica deliberada con caveat de efecto acotado (Ericsson;
%   Charness; Macnamara), efectuación (Sarasvathy), formación de hábitos en 66 días (Lally),
%   intenciones de implementación (Gollwitzer y Sheeran), acumulación conductual de riqueza
%   (Thaler y Benartzi).
%   Sin secciones de refutación; únicamente ciencia validada y accionable.
% Enlazar fuertemente con \cref{ch:decisions-under-uncertainty} (sesgos) y
%   \cref{ch:organizational-behavior} (motivación/SDT) --- NO reescribir; usar \cref.

La literatura sobre psicología del emprendedor es extensa, controvertida y, con frecuencia,
escrita al revés: comienza con los fundadores exitosos y reconstruye hacia atrás lo que afirman
haber creído, hecho o sentido. Este capítulo opera en sentido contrario. Parte de metaanálisis
revisados por pares y estudios de campo prerre\-gis\-tra\-dos, extrae los mecanismos y tamaños
del efecto que superan el escrutinio riguroso, y los mapea sobre las decisiones que un fundador
enfrenta en la realidad. Los hallazgos son más modestos de lo que sugieren las portadas de los
\emph{bestsellers}, pero son más útiles precisamente porque son accionables: metas específicas,
confianza anclada al dominio, rasgos adaptados a la tarea emprendedora, experiencia construida
mediante retroalimentación estructurada, decisiones calibradas a la pérdida asequible, y
comportamiento codificado en sistemas en lugar de depender de la fuerza de voluntad. Estos son
los constructos que cuentan con replicación. Y, como demuestran las secciones siguientes,
están estrechamente vinculados a las métricas operativas de la empresa SaaS que atraviesa
todos los ejemplos trabajados del libro.

\section{Autoeficacia y establecimiento de metas}\label{sec:self-efficacy-goals}
% complexity: 2

¿Qué separa a un fundador que intenta resolver un problema difícil de uno que simplemente lo
discute? La respuesta próxima no es el talento, el capital ni la suerte --- es la creencia de
que un curso de acción específico, ejecutado por ellos, producirá un resultado que pueden
influir. Albert Bandura formalizó este constructo como \emph{autoeficacia}: el juicio que un
individuo hace sobre su capacidad para organizar y ejecutar los cursos de acción que una
situación específica requiere \parencite{bandura1997self}. Dos características distinguen la
autoeficacia de conceptos más vagos como la confianza general o el optimismo. En primer lugar,
es específica al dominio: un fundador puede tener autoeficacia alta para la iteración de
producto y baja para ventas empresariales en la misma semana. En segundo lugar, predice el
comportamiento mediante mecanismos y no solo por correlación --- las personas con alta
autoeficacia en una tarea abordan los obstáculos como problemas a dominar, sostienen el
esfuerzo por más tiempo cuando encuentran resistencia y se recuperan más rápido tras los reveses.

La evidencia metaanalítica sobre este mecanismo es inusualmente sólida. Stajkovic y Luthans
sintetizaron \num{114} estudios independientes que abarcan \num{21616} participantes y hallaron
una correlación promedio ponderada entre la autoeficacia específica a la tarea y el rendimiento
laboral de $r \approx .38$ % fuente: stajkovic1998self, Psychological Bulletin 1998, k=114, N=21,616
\parencite{stajkovic1998self}. En términos prácticos, pasar de baja a alta autoeficacia en un
dominio determinado se asocia con un incremento de aproximadamente \qty{28}{\percent} en la
producción, que opera a través de tres canales: las personas con alta eficacia seleccionan metas
más difíciles y de mayor rendimiento, aplican mayor esfuerzo inicial al comenzar una tarea y
persisten por más tiempo cuando encuentran resistencia o rechazo del mercado. El canal importa
tanto como la correlación: la autoeficacia no mejora el rendimiento haciendo que los fundadores
se sientan mejor; lo mejora cambiando las metas que se fijan y el esfuerzo que sostienen. La
motivación de logro --- el impulso intrínseco de superar desafíos y alcanzar estándares objetivos
elevados --- amplifica esto aún más. El metaanálisis de Collins, Hanges y Locke demostró que la
motivación de logro está significativamente correlacionada tanto con la elección de una carrera
emprendedora como con el desempeño emprendedor, con una diferenciación particularmente marcada
al comparar fundadores de alto rendimiento con sus pares de menor desempeño % fuente: collins2004achievement, Human Performance 2004
\parencite{collins2004achievement}.

La autoeficacia responde a \emph{si} actuar. El establecimiento de metas responde a \emph{hacia
dónde} apuntar. La síntesis metaanalítica de Locke y Latham, construida a lo largo de treinta y
cinco años de estudios de campo y de laboratorio, llegó a una conclusión engañosamente simple:
las metas específicas y desafiantes producen sistemáticamente un rendimiento superior al de las
exhortaciones vagas de «haz tu mejor esfuerzo», con tamaños del efecto ($d$) que oscilan entre
\num{0.52} y \num{0.82} en los ensayos sintetizados \parencite{locke2002building}. El mecanismo
opera a través de cuatro canales: las metas dirigen la atención hacia las actividades relevantes
y la alejan de las irrelevantes; las metas elevadas incrementan la energía aplicada a la tarea;
las metas difíciles sostienen el esfuerzo por más tiempo porque la brecha entre el estado actual
y el objetivo mantiene activa la motivación; y las metas desencadenan el descubrimiento de
estrategias --- cuando los enfoques existentes son insuficientes, la meta crea presión para buscar
mejores alternativas.

Dos condiciones delimitan el efecto. En primer lugar, el compromiso con la meta: si un objetivo
se siente impuesto políticamente o genuinamente imposible, el mecanismo de cuatro canales nunca
se activa. En segundo lugar, la complejidad de la tarea: para tareas rutinarias, la regla
«específico y desafiante» es robusta; para problemas genuinamente novedosos --- situación común
en las empresas en etapa temprana --- los objetivos de rendimiento rígidos pueden producir un
comportamiento desordenado y poco sistemático. En esas condiciones, Locke y Latham recomiendan
sustituir una \emph{meta de aprendizaje} («descubrir los tres factores que impulsan la
conversión de prueba a pago») por una meta de rendimiento («alcanzar \qty{25}{\percent} de
conversión este mes»). La distinción importa para los fundadores porque su entorno mezcla
ambas: la ejecución sobre procesos conocidos merece objetivos de rendimiento exigentes; la
exploración genuina merece el encuadre de una meta de aprendizaje.

\begin{example}[Fundador SaaS que establece una estructura de hitos]
La fundadora de la empresa SaaS con \money{340000} en ingresos recurrentes mensuales decide
mejorar la tasa de conversión de prueba a pago, que actualmente se encuentra en
\qty{25}{\percent}. Reconoce que la dinámica de canales y los problemas de atribución ya están
bien comprendidos a partir del trabajo realizado en \cref{ch:attribution-and-measurement}; la
pregunta de ejecución está lo suficientemente resuelta como para que un objetivo de rendimiento
sea apropiado.

Siguiendo a Locke y Latham, establece una meta específica y desafiante: alcanzar
\qty{30}{\percent} de conversión de prueba a pago en noventa días. Luego la descompone en hitos
semanales --- instrumentar el flujo de incorporación (\emph{onboarding}) antes del día siete;
ejecutar la primera prueba A/B del correo de activación antes del día catorce; evaluar las
diferencias de retención por cohorte antes del día treinta. Cada hito tiene un resultado medible
y una fecha límite, preservando las funciones directiva y energizante de la teoría del
establecimiento de metas. Dado que el proceso de conversión en sí es bien conocido (no novedoso),
mantiene el encuadre de rendimiento a lo largo de todo el proceso. Complementa la meta con una
verificación deliberada de autoeficacia: habiendo lanzado tres experimentos A/B previos, dispone
de evidencia específica al dominio de que puede ejecutar este tipo de proyecto, lo que aumenta
la probabilidad de que sostenga el esfuerzo cuando el primer experimento decepcione.

La implicación es clara: el establecimiento de metas funciona precisamente porque la estructura
del objetivo --- específico, con fecha, descompuesto --- hace imposible sustituir la intención
vaga por la acción concreta. El aumento de rendimiento de \qty{28}{\percent} asociado con la alta
autoeficacia específica a la tarea se combina con el efecto del establecimiento de metas cuando
ambos están presentes; ninguno opera bien en ausencia del otro.
\end{example}

\begin{admonition}{Nota}
La autoeficacia se construye, no se hereda. Bandura identifica cuatro fuentes: experiencias de
dominio (éxitos previos en el mismo dominio), modelado vicario (observar a otros comparables
tener éxito), persuasión social (retroalimentación creíble de personas que han observado el
desempeño), y estado fisiológico (interpretar la ansiedad como activación y no como
incompetencia). La implicación práctica para los fundadores es que la autoeficacia en un nuevo
dominio --- ventas empresariales, modelado financiero, hablar en público --- se incrementa de
manera más confiable acumulando pequeños logros específicos al dominio, no mediante insumos
motivacionales generales.
\end{admonition}

La literatura de motivación intrínseca, revisada en \cref{sec:motivation}, converge con el
modelo de Bandura: la autonomía y la competencia son las dos necesidades de la teoría de la
autodeterminación más directamente activadas por un sistema de metas bien estructurado. La meta
otorga a la tarea una estructura clara (competencia), mientras que la elección del objetivo y
el método por parte del fundador preserva la apropiación (autonomía). \textcite{bandura1997self},
\textcite{stajkovic1998self} y \textcite{locke2002building} proveen la base empírica; la
pregunta práctica es siempre si la tarea está suficientemente resuelta para una meta de
rendimiento o si es lo bastante novedosa como para requerir primero una meta de aprendizaje.

\section{Rasgos y capital humano que predicen el desempeño de la empresa}\label{sec:traits-human-capital}
% complexity: 3

Si la autoeficacia y el compromiso con la meta explican cómo los fundadores sostienen el
esfuerzo, ¿qué características personales predicen si tendrán éxito en la tarea de dirigir un
negocio? Durante décadas, la investigación en emprendimiento produjo respuestas inconsistentes
porque formulaba la pregunta equivocada: si los rasgos de personalidad amplios --- extraversión,
amabilidad, apertura a la experiencia --- predecían los resultados de la empresa. En su mayoría,
no lo hacen. El avance llegó al reformular la pregunta: ¿qué rasgos están lógicamente
\emph{adaptados} a las demandas específicas de la tarea de dirigir un negocio, y qué dimensiones
de capital humano se \emph{despliegan} realmente en lugar de simplemente acumularse?

El metaanálisis de Rauch y Frese proporcionó la cuantificación definitiva. Al analizar variables
de personalidad en un gran conjunto de muestras independientes, hallaron un contraste marcado
entre rasgos adaptados versus no adaptados a la tarea emprendedora. Los rasgos adaptados ---
necesidad de logro, tolerancia al estrés, personalidad proactiva, necesidad de autonomía,
innovación y autoeficacia generalizada --- produjeron correlaciones con la creación de negocios
de $r = .247$ ($K = 47$, $N = \num{13280}$) y con el éxito empresarial de $r = .250$
($K = 42$, $N = \num{5607}$). Los rasgos no adaptados mostraron una validez predictiva
despreciable para el éxito ($r = .028$, $K = 13$, $N = \num{2777}$). % fuente: rauch2007person, European Journal of Work & Org Psychology 2007
La implicación práctica es que la extraversión general o la responsabilidad como rasgos
abstractos son predictores débiles; la tolerancia al estrés y la iniciativa proactiva --- rasgos
estructuralmente adaptados a las demandas de un entorno volátil y de alto riesgo --- sí
representan una señal real \parencite{rauch2007person}.

Los hallazgos sobre capital humano corren en paralelo. Los supuestos empresariales
convencionales confunden dos cosas que son significativamente distintas: las \emph{inversiones}
que un individuo ha realizado en capital humano (años de educación formal, trayectoria gerencial
general, un MBA) versus los \emph{resultados} que esas inversiones pueden o no haber producido
(conocimiento y habilidades específicos, demostrables y relacionados con la tarea). Unger, Rauch,
Frese y Rosenbusch sintetizaron \num{70} muestras independientes con un total de
$N = \num{24733}$ empresas y encontraron una relación general entre el capital humano general y
el éxito empresarial de $r_c = .098$ --- pequeña en términos absolutos
\parencite{unger2011human}. % fuente: unger2011human, Journal of Business Venturing 2011, N=24,733
Pero el agregado oculta el hallazgo operativamente crítico: los \emph{resultados} de capital
humano predijeron el éxito a tasas sustancialmente más altas que las \emph{inversiones} en
capital humano, y el capital humano relacionado con la tarea superó ampliamente al capital de
baja relación con la tarea. Una década de experiencia gerencial general es un predictor débil;
la capacidad específica de gestionar un proceso de ventas empresariales de software de ciclo
completo o diseñar una arquitectura de datos escalable es un predictor sólido. El motivo es
directo: el éxito de una empresa está determinado por la calidad de las decisiones tomadas en
contextos específicos y de alto riesgo, y las credenciales genéricas señalizan poco sobre el
desempeño en esos contextos.

Para los fundadores, estos dos hallazgos --- rasgos adaptados y capital humano orientado a
resultados --- tienen implicaciones concretas para la autoevaluación y la contratación temprana.
El rasgo que merece atención no es cuán extrovertido o agradable es un cofundador, sino si
demuestra la tolerancia al estrés y la orientación proactiva que Rauch y Frese identificaron
como adaptadas a la tarea. El criterio de contratación que merece atención no son las
credenciales ni los años de experiencia, sino la capacidad de demostrar habilidades relevantes
para la tarea en una evaluación estructurada.

\begin{example}[Auditoría de rasgos y capital humano para una decisión de contratación en la SaaS]
La fundadora de la SaaS, con \money{340000} en ingresos recurrentes mensuales, \money{3000000}
en capital invertido y una tasa de corte del \qty{11}{\percent}, evalúa a dos candidatos para el
puesto de jefa de ventas. El candidato A tiene un MBA de una institución reconocida y pasó ocho
años en la gestión de cuentas corporativas. El candidato B tiene un historial general más breve,
pero ha cerrado demostrablemente seis contratos empresariales de SaaS en los últimos dos años
con valores de contrato promedio superiores a \money{60000}, y maneja sin evasión el ejercicio
de \emph{roleplay} de llamada de ventas estructurada durante la entrevista.

El \emph{framework} de Unger et al.\ la dirige a ponderar el capital humano orientado a
resultados del candidato B --- habilidad específica, desplegada y relacionada con la tarea ---
por encima de las credenciales de inversión del candidato A. El \emph{framework} de Rauch y
Frese refuerza esto: si el candidato B también puntúa más alto en tolerancia al estrés
observable e iniciativa proactiva de tareas durante el proceso de entrevista (respondiendo al
rechazo con un enfoque revisado en lugar de retirarse), esos rasgos adaptados multiplican la
ventaja de capital humano. El margen de contribución en juego --- \money{21} al mes por cliente
convertido por la contratación de ventas --- le proporciona un marco cuantitativo sobre el valor
de la decisión de contratación. Cien clientes adicionales por año con un margen anual de
\money{252} cada uno, descontados a la tasa de corte del \qty{11}{\percent} como una perpetuidad
creciente con crecimiento del \qty{3}{\percent}, representa aproximadamente \money{315000} de
valor presente. Una decisión de contratación calibrada sobre evidencia de resultados y no sobre
señalización de credenciales es directamente trazable a ese número.

La implicación: las evaluaciones de rasgos adaptados y capital orientado a resultados no son
complementos «blandos» de un proceso de contratación; son los predictores cuantitativamente
mejor respaldados de si la contratación producirá los resultados de métricas que la empresa
necesita.
\end{example}

\begin{admonition}{Nota}
La literatura sobre sesgos cognitivos es directamente relevante en este dominio. La
sobreponderación de credenciales y experiencia general en la contratación refleja tanto la
heurística de representatividad («esta persona se parece a las personas exitosas que conozco»)
como el sesgo de disponibilidad (los éxitos de líderes con muchas credenciales son más visibles
que sus fracasos). Los datos de Rauch y Frese son un antídoto empírico: la pregunta relevante
no es si el candidato coincide con el prototipo de un líder exitoso, sino si sus rasgos y
habilidades demostradas están adaptados a las demandas específicas de la tarea del puesto. El
andamiaje de teoría de decisiones en \cref{ch:decisions-under-uncertainty} desarrolla esto
directamente.
\end{admonition}

Ambos conjuntos de hallazgos se conectan con la columna vertebral de creación de valor: los
resultados de capital humano no son un dotación fija sino el producto acumulado de experiencia
deliberada y rica en retroalimentación, que es el tema de la siguiente sección.
\textcite{rauch2007person} y \textcite{unger2011human} proveen la base metaanalítica;
\textcite{collins2004achievement} conecta el hilo de motivación de logro de
\cref{sec:self-efficacy-goals} con el panorama más amplio de la personalidad.

\section{Cómo se construye la experiencia y cómo deciden los fundadores}\label{sec:practice-effectuation}
% complexity: 3

Si los rasgos adaptados y el capital humano orientado a resultados predicen el desempeño de la
empresa, la siguiente pregunta es práctica: ¿cómo construye un fundador esos resultados, y cómo
debe desplegarse la experiencia resultante en el entorno genuinamente incierto de una empresa en
etapa temprana? La respuesta proviene de dos cuerpos de investigación que generalmente se
enseñan por separado --- ciencia de la experiencia experta y cognición emprendedora --- pero que
son lógicamente secuenciales.

Ericsson, Krampe y Tesch-R{\"o}mer introdujeron la práctica deliberada para distinguir el
rendimiento experto de la mera experiencia acumulada \parencite{ericsson1993role}. Las
propiedades clave son la especificidad, la retroalimentación inmediata y la dificultad
progresivamente creciente. La práctica debe apuntar a una debilidad bien definida; la
retroalimentación sobre el rendimiento debe ser rápida y precisa; y el desafío debe calibrarse
justo por encima del nivel de competencia actual. El trabajo que cae dentro de la competencia
establecida la consolida, pero no la extiende. Ericsson pasó la última parte de su carrera
frustrado por la reducción popular de sus hallazgos a «diez mil horas de cualquier cosa», porque
la investigación muestra que diez mil horas de experiencia no estructurada produce un buen
amateur, no un experto.

La evidencia de campo sobre la experiencia en ajedrez confirma el mecanismo. Charness, Tuffiash,
Krampe, Reingold y Vasyukova demostraron que el estudio solitario y serio --- práctica deliberada
en el sentido técnico --- era el predictor más sólido del nivel de habilidad entre los jugadores,
explicando gran parte de la varianza en las calificaciones Elo que la experiencia acumulada sola
no podía explicar \parencite{charness2005practice}. % fuente: charness2005practice, Applied Cognitive Psychology 2005
El dominio del ajedrez es instructivo precisamente porque tiene reglas estables, métricas
objetivas claras y ciclos de retroalimentación rápidos --- las condiciones bajo las cuales la
práctica deliberada es más poderosa.

Los practicantes deben aplicar una advertencia crítica al trasladar esto al entorno empresarial.
El metaanálisis de Macnamara, Hambrick y Oswald cuantificó la proporción de varianza del
rendimiento explicada por la práctica deliberada en múltiples dominios. % fuente: macnamara2014practice, Psychological Science 2014
En juegos con reglas estables (incluido el ajedrez), la práctica deliberada explicó el
\qty{26}{\percent} de la varianza en el rendimiento. En música, explicó el \qty{21}{\percent}.
En dominios profesionales menos estructurados con retroalimentación lenta y ruidosa, la cifra
cae sustancialmente \parencite{macnamara2014practice}. Para los fundadores, la implicación es
calibrada y no desalentadora: la práctica deliberada es muy eficaz para sub-habilidades
aisladas y bien definidas --- manejo de objeciones en ventas empresariales, modelado financiero
en una hoja de cálculo, redacción de un guión de entrevista de descubrimiento de clientes ---
pero no es el mecanismo principal para «construir una empresa» como actividad holística y
ambigua. El hallazgo de capital humano de Unger et al.\ (\cref{sec:traits-human-capital})
refuerza esto: el capital humano que predice el éxito empresarial está relacionado con la tarea
y es desplegable, no genérico.

El entorno emprendedor hace difícil construir una práctica deliberada genuina a nivel de la
empresa porque los ciclos de retroalimentación son lentos, ruidosos y costosos. Un cirujano sabe
en minutos si un procedimiento tuvo éxito; un emprendedor puede esperar doce meses para saber si
una decisión de producto fue correcta. La respuesta práctica es crear estructuras de
retroalimentación artificiales: las pruebas A/B rigurosas proveen la señal rápida y objetiva que
un entrenador provee en un dominio de experiencia tradicional; el modelado financiero expone los
supuestos a la falsificación antes de que lo haga el mercado; las relaciones de asesoramiento
con operadores experimentados funcionan como el \emph{coach} experto que el entorno no provee
automáticamente.

La toma de decisiones experta en el emprendimiento tiene una segunda característica distintiva,
identificada por Sarasvathy mediante un estudio pionero de treinta fundadores expertos ---
emprendedores con al menos quince años de experiencia que habían llevado al menos una empresa a
cotizar en bolsa \parencite{sarasvathy2001causation}. Ante el mismo caso de negocio, los
fundadores expertos no razonaban de la manera que prescriben los libros de texto de MBA. No
partían de un mercado objetivo, proyectaban retornos y movilizaban recursos hacia el objetivo.
En cambio, partían de sus medios disponibles --- quiénes son, qué saben, a quién conocen --- y
permitían que los objetivos emergieran de lo que esos medios podían plausiblemente producir. Las
pérdidas estaban acotadas por un umbral de «pérdida asequible»: no el retorno esperado máximo,
sino el máximo que estaban dispuestos a arriesgar. Sarasvathy denominó a este patrón
\emph{efectuación}, en contraste con la lógica de \emph{causación} de la planificación
estratégica estándar.

La distinción importa porque no es simplemente una preferencia estilística. Los fundadores
expertos usan la efectuación porque está mejor adaptada a entornos de incertidumbre genuina,
donde las distribuciones de probabilidad requeridas por el razonamiento de valor esperado no
existen. En ese régimen, controlar el riesgo a la baja y aprovechar las contingencias a medida
que surgen supera al intento de predecir y optimizar. El enfoque causal --- investigación de
mercado, retorno proyectado, adquisición de recursos --- es apropiado cuando el futuro es
suficientemente predecible y el mercado es suficientemente estable. El emprendimiento en etapa
temprana raramente satisface ninguna de las dos condiciones. Esto se conecta directamente con el
andamiaje de teoría de decisiones en \cref{ch:decisions-under-uncertainty}: la efectuación es,
en esencia, la implementación práctica de la toma de decisiones robusta bajo incertidumbre
profunda.

\begin{example}[Encuadre efectual vs.\ causal para la SaaS]
La fundadora de la SaaS, actualmente con \money{340000} en ingresos recurrentes mensuales y
una tasa de retención a doce meses del \qty{92.5}{\percent}, está considerando ingresar a un
nuevo vertical --- despachos de abogados pequeños. Un tomador de decisiones causal comenzaría
estimando el mercado total direccionable, proyectando una tasa de penetración, modelando el
margen de contribución del nuevo vertical (\money{21} al mes por cliente, según la estructura de
costos de la empresa), calculando el valor presente neto de la expansión usando la tasa de corte
del \qty{11}{\percent}, y adquiriendo recursos para ejecutar si el NPV es positivo.

Una fundadora efectual hace preguntas iniciales diferentes: ¿Conocemos a socios en despachos de
abogados pequeños que lo intentarían de forma gratuita y nos darían retroalimentación honesta?
¿Podemos configurar el producto existente para el vertical dentro del \emph{sprint} de
ingeniería actual, o requiere una base de código separada? ¿Cuál es el tiempo y presupuesto
máximos que estamos dispuestos a perder si esto resulta ser un callejón sin salida --- y ese
monto es genuinamente asequible contra la \emph{runway} actual? No se niega a modelar retornos,
pero trata el modelo como uno de varios insumos y pondera la evidencia accesible de medios ---
relaciones, capacidad existente, exposición asequible --- con mayor peso que la evidencia
proyectada de un pronóstico de mercado.

La implicación: los dos \emph{frameworks} no compiten en todas las situaciones. El enfoque
causal se vuelve apropiado una vez que la respuesta del mercado está mejor comprendida; la
efectuación es la postura correcta cuando la incertidumbre es genuina y no puede resolverse
mediante más investigación.
\end{example}

\begin{admonition}{Nota}
La práctica deliberada y la efectuación no están en tensión; son secuenciales. La efectuación
es la lógica de decisión para entrar en una situación incierta con riesgo acotado. La práctica
deliberada es lo que construye la competencia de dominio que hace que la efectuación futura sea
más rica --- porque los medios del fundador (lo que sabe, lo que puede hacer) se expanden con
cada iteración rica en retroalimentación. Los emprendedores expertos del estudio de Sarasvathy
tenían extensas historias de práctica deliberada en sus dominios; la efectuación era la forma
en que aprovechaban esa experiencia en situaciones nuevas, no un sustituto de haberla
construido.
\end{admonition}

\section{Hábitos y acumulación conductual de riqueza}\label{sec:habits-wealth}
% complexity: 2

El enfoque del alto rendimiento como producto de la fuerza de voluntad y la motivación es uno de
los supuestos más consistentemente erróneos de la psicología popular. La fuerza de voluntad se
agota; la motivación fluctúa; ninguna es una base confiable para el cambio de comportamiento a
largo plazo que el crecimiento compuesto --- personal o financiero --- requiere. La ciencia
conductual de la formación de hábitos y la arquitectura de elección ofrece una alternativa más
duradera: diseñar el entorno de modo que el comportamiento deseado sea el camino de menor
resistencia, y el comportamiento indeseado requiera una decisión activa para acceder.

El modelo científico fundamental para entender cómo los hábitos interactúan con las metas es
provisto por Wood y Neal, quienes demostraron que los hábitos emergen a través del aprendizaje
gradual de asociaciones entre respuestas y características específicas de los contextos de
ejecución --- hora del día, ubicación física o secuencias de acciones precedentes
\parencite{wood2007habits}. Una vez que un hábito se forma, la percepción de la señal contextual
desencadena la respuesta conductual asociada sin la mediación de un estado de meta activo. Las
metas son necesarias para motivar las repeticiones iniciales que establecen el hábito, pero una
vez que se alcanza la automaticidad, el comportamiento se vuelve sustancialmente resistente a
los cambios en el estado motivacional del individuo. Para los fundadores, esta automaticidad
preserva la función ejecutiva para la resolución de problemas complejos --- el recurso cognitivo
más en riesgo durante las fases de alto estrés de una empresa.

La pregunta práctica que los libros populares sobre hábitos no responden de manera rigurosa es
cuánto tiempo toma este proceso de formación. Charles Duhigg, James Clear y B.\ J.\ Fogg han
descrito la arquitectura de hábitos de forma convincente, y sus \emph{frameworks} son guías de
diseño útiles a nivel estructural \parencite{duhigg2012power,clear2018atomic,fogg2020tiny}. Pero
la afirmación del «hábito de 21 días» arraigada en la cultura popular no está respaldada
empíricamente. Lally, van Jaarsveld, Potts y Wardle hicieron seguimiento a noventa y seis
voluntarios durante ochenta y cuatro días mientras intentaban adoptar un nuevo comportamiento de
salud vinculado a una señal diaria, mapeando la automaticidad mediante el Índice de Hábitos de
Autoinforme \parencite{lally2010habit}. % fuente: lally2010habit, European Journal of Social Psychology 2010, N=96, 84 días
La curva de automaticidad fue asintótica: grandes ganancias en las repeticiones iniciales,
rendimientos decrecientes posteriormente, hasta alcanzar una meseta en un máximo individual.
El tiempo promedio requerido para alcanzar el \qty{95}{\percent} de la asíntota de automaticidad
fue de \textbf{66 días}, con varianza individual que oscilaba entre 18 y 254 días dependiendo
de la complejidad del comportamiento. Crucialmente, omitir una única oportunidad de ejecutar el
comportamiento no afectó materialmente el proceso de formación ni alteró la asíntota final ---
lo que significa que las rachas ininterrumpidas rígidas no son neurológicamente necesarias, pero
un horizonte sostenido de más de sesenta días sí lo es.

Dado que los hábitos tardan una mediana de dos meses en formarse, los fundadores requieren un
puente confiable entre la intención y la automaticidad. Ese puente es la intención de
implementación. Gollwitzer y Sheeran analizaron este mecanismo en noventa y cuatro pruebas
independientes y encontraron un efecto robusto de magnitud media a grande sobre el logro de
metas ($d \approx .65$) \parencite{gollwitzer2006implementation}.
% fuente: gollwitzer2006implementation, Advances in Experimental Social Psychology 2006, 94 pruebas, d≈.65
Una intención de implementación es un plan explícito «si-entonces» que especifica cuándo y
dónde se iniciará un comportamiento dirigido a una meta: «Si son las 8:00 AM de un día de
semana, entonces abriré el \emph{dashboard} de MRR antes que cualquier otra aplicación.» Este
formato transfiere el desencadenante de la acción desde la fuerza de voluntad interna hacia una
señal ambiental externa, promoviendo el inicio rápido del comportamiento dirigido a la meta,
protegiendo la persecución en curso del objetivo de demandas competidoras, y reduciendo la carga
cognitiva necesaria para comenzar. El mecanismo clave es la especificidad: el formato «si-entonces»
pre-compromete la respuesta a una señal contextual, de modo que cuando la señal aparece, el
comportamiento se ejecuta sin requerir una decisión.

La aplicación de finanzas conductuales extiende la misma lógica a las decisiones financieras
personales. Los emprendedores enfrentan una versión aguda del problema estándar de autocontrol:
ingresos irregulares, fuertes preferencias sesgadas hacia el presente y la tentación constante
de reasignar los ahorros personales al negocio. El programa «Save More Tomorrow» de Thaler y
Benartzi provee la estrategia contra-intuitiva basada en evidencia \parencite{thaler2004save}.
El \emph{insight} clave de diseño del programa es que la aversión a la pérdida --- la asimetría
por la cual las pérdidas pesan más que las ganancias equivalentes en la experiencia subjetiva ---
hace dolorosos los recortes inmediatos al ingreso disponible, incluso cuando el beneficio a largo
plazo es obvio. «Save More Tomorrow» sortea esto pidiendo a los participantes que se pre-comprometan
a destinar una fracción fija de sus \emph{futuros} aumentos de sueldo al ahorro. Dado que la
deducción coincide con un incremento en el salario bruto, el ingreso disponible nunca cae; no
se activa la aversión a la pérdida. Thaler y Benartzi hicieron seguimiento a los participantes
durante cuarenta meses y observaron que las tasas promedio de ahorro aumentaron del
\qty{3.5}{\percent} al \qty{13.6}{\percent} --- una mejora de casi cuatro veces lograda sin
fuerza de voluntad, solo mediante arquitectura. El mecanismo es idéntico al de la literatura
sobre hábitos e intenciones de implementación: eliminar la decisión del ámbito de la deliberación
en el momento; codificarla de antemano como un compromiso automático.

\begin{example}[Intenciones de implementación y ahorro automatizado para la fundadora SaaS]
La fundadora de la SaaS quiere establecer dos rutinas confiables: una revisión financiera semanal
y un protocolo personal de acumulación de riqueza. Aplica el hallazgo de Lally et al.\ de que
la automaticidad del hábito requiere un horizonte mediano de sesenta y seis días, y el hallazgo
de Gollwitzer y Sheeran de que un formato explícito «si-entonces» ($d \approx .65$) tiende el
puente mientras el hábito se forma.

Para la revisión financiera, escribe explícitamente la intención de implementación: «Si es el
viernes por la mañana y mi calendario muestra las 9:00 AM, entonces abriré el \emph{dashboard}
de MRR, verificaré la tasa de \emph{burn} contra la base de \money{3000000} de capital invertido,
y registraré una decisión o hipótesis actualizada antes de cerrarlo.» La señal (viernes 9:00 AM)
es estable y externa. La rutina (secuencia de tres elementos) es lo suficientemente específica
como para ejecutarse sin mayor deliberación. Después de diez semanas --- dentro del rango de
Lally et al.\ --- la revisión se ha convertido en algo activado por el contexto y no dependiente
de la motivación: ocurre tanto si se siente comprometida con ello como si no.

En el lado de la acumulación de riqueza, aplica directamente la arquitectura de Thaler y
Benartzi: cada vez que el negocio realiza una distribución trimestral, el \qty{20}{\percent} se
redirige automáticamente a una cuenta de inversión personal antes de llegar a su cuenta operativa.
La instrucción está fijada en el sistema de órdenes permanentes del banco y requiere una acción
deliberada para cancelarse. El desencadenante es un evento de entrada de efectivo, no una
decisión de gasto, por lo que la aversión a la pérdida no se activa. El equivalente a
pre-comprometerse con futuros aumentos de sueldo es pre-comprometerse con futuras fracciones de
distribución: el sacrificio es prospectivo, no inmediato.

La implicación es general: tanto el hábito profesional como la regla de ahorro funcionan porque
son arquitectura, no determinación. El formato de intención de implementación maneja el puente
de cuarenta a sesenta días antes de la automaticidad; el bloqueo de arquitectura de elección
maneja el valor predeterminado a largo plazo. En el momento en que cualquiera de los dos requiere
que la fundadora decida en tiempo real, bajo presiones competidoras, se vuelven poco confiables.
\end{example}

\begin{admonition}{Nota}
La necesidad de autonomía identificada en la teoría de la autodeterminación (\cref{sec:motivation})
no queda socavada por la automatización del hábito --- se expresa en el momento de diseñar el
sistema. La fundadora que se pre-compromete con un esquema de ahorro y escribe una intención de
implementación para una revisión financiera semanal ha ejercido alta autonomía al construir la
arquitectura de elección; lo que se automatiza después es la ejecución, no la elección. La
distinción importa porque los fundadores a veces resisten los enfoques sistemáticos como
restricciones a su agencia, cuando en realidad diseñar sistemas eficaces es uno de los ejercicios
de mayor apalancamiento de agencia disponibles para ellos.
\end{admonition}

El argumento del capítulo cierra donde comenzó: la psicología emprendedora que sobrevive el
contacto con la replicación no es un perfil de personalidad ni un sistema de creencias. Es un
conjunto de prácticas --- metas estructuradas acompañadas de confianza específica al dominio
construida mediante evidencia, rasgos y capital humano adaptados a las demandas de la tarea,
lógica de decisión adaptada a la incertidumbre genuina, y comportamiento codificado en sistemas
en lugar de resolverse en tiempo real. Estas son las herramientas. \textcite{bandura1997self},
\textcite{stajkovic1998self}, \textcite{locke2002building}, \textcite{rauch2007person},
\textcite{unger2011human}, \textcite{ericsson1993role}, \textcite{sarasvathy2001causation},
\textcite{lally2010habit}, \textcite{gollwitzer2006implementation} y \textcite{thaler2004save}
proveen la base empírica; los \emph{frameworks} en \cref{ch:decisions-under-uncertainty} y
\cref{sec:motivation} los ubican dentro de la ciencia de toma de decisiones y motivación más
amplia que el libro ya ha desarrollado.
